% F6 fig:order-ladder (Chapter 4, text width). Relative error against the step for Central<2> to Central<10>,
% exp(x) at x = 1 (DS2, data/step-error-exp.csv). The truncation branches steepen with P; the rounding branches
% nearly coincide (S from 1 to 2.3), and the minima slide up the common rounding line eps / h (muted, dashed),
% because a higher order lets the step grow. The five curves have five colours, ordered dark to light with P so
% that they also read in greyscale, and each label takes its curve's colour (the author's release decision: the
% grey ladder was hard to read). The float and caption stay in the chapter.
\definecolor{ladderTwo}{HTML}{2B2B2B}
\definecolor{ladderFour}{HTML}{0B5FAE}
\definecolor{ladderSix}{HTML}{1B8F8A}
\definecolor{ladderEight}{HTML}{6AA84F}
\definecolor{ladderTen}{HTML}{C79A1E}
\begin{tikzpicture}
\begin{axis}[eqtufte, width=95mm, height=45mm,
  xmode=log, ymode=log, log basis x=10, log basis y=10,
  xmin=1e-13, xmax=1, ymin=1e-16, ymax=1,
  xtick={1e-12,1e-10,1e-8,1e-6,1e-4,1e-2,1}, ytick={1e-16,1e-12,1e-8,1e-4,1},
  xlabel={step $h$}, ylabel={relative error}, unbounded coords=jump]
  % the common rounding line eps / h (relative, since f = f' for exp): a guide, not data
  \addplot[black!45, dashed, thin, domain=1e-13:1e-3, samples=2] {2^(-52) / x};
  \node[black!55, font=\scriptsize, anchor=south west, inner sep=1pt] at (axis cs:3e-9, 1.5e-7) {$\epsmach / h$};
  \addplot[ladderTwo] table[col sep=comma, comment chars={\#}, x=h, y=central2] {data/step-error-exp.csv};
  \addplot[ladderFour] table[col sep=comma, comment chars={\#}, x=h, y=central4] {data/step-error-exp.csv};
  \addplot[ladderSix] table[col sep=comma, comment chars={\#}, x=h, y=central6] {data/step-error-exp.csv};
  \addplot[ladderEight] table[col sep=comma, comment chars={\#}, x=h, y=central8] {data/step-error-exp.csv};
  \addplot[ladderTen] table[col sep=comma, comment chars={\#}, x=h, y=central10] {data/step-error-exp.csv};
  % labels on the truncation branches, left of each curve in its colour (the curve is lower to the left)
  \node[ladderTwo, font=\scriptsize, anchor=east, inner sep=1pt] at (axis cs:0.17, 5e-3) {\code{Central<2>}};
  \node[ladderFour, font=\scriptsize, anchor=east, inner sep=1pt] at (axis cs:0.17, 3e-5) {\code{Central<4>}};
  \node[ladderSix, font=\scriptsize, anchor=east, inner sep=1pt] at (axis cs:0.17, 2e-7) {\code{Central<6>}};
  \node[ladderEight, font=\scriptsize, anchor=east, inner sep=1pt] at (axis cs:0.17, 1.6e-9) {\code{Central<8>}};
  \node[ladderTen, font=\scriptsize, anchor=east, inner sep=1pt] at (axis cs:0.17, 1.2e-11) {\code{Central<10>}};
\end{axis}
\end{tikzpicture}
